% Generated by scripts/assemble.py; corrections belong in source/corrections.json.
\examhead{161}{华南师范大学 2026 年数学分析试卷}{153}
% p154-b20; source PDF page 154
\par\Needspace{5\baselineskip}
\blockanchor{p154-b20}
\majorquestion{一}求极限. 每题 10 分, 共 50 分.\par
% p154-b21; source PDF page 154
\blockanchor{p154-b21}
\question{1}$\lim_{n\to \infty}\sqrt{n}\left(\sqrt{n^2 + 1} -\sqrt{n + 1}\right).$\par
% p154-b22; source PDF page 154
\blockanchor{p154-b22}
\question{2}$\lim_{n\to\infty}n(\sqrt[n]{a}-1)(a>0).$\par
% p154-b23; source PDF page 154
\blockanchor{p154-b23}
\question{3}$\lim_{x\to0}\frac{\sqrt{\cos x}-\sqrt{1+x\arctan x}}{2x^{2}}$ .\par
% p154-b24; source PDF page 154
\blockanchor{p154-b24}
\question{4}$\lim_{x\to+\infty}e^{-x}\left(1+\frac{1}{x}\right)^{x^{2}}$\par
% p154-b25; source PDF page 154
\blockanchor{p154-b25}
\question{5}$\lim_{(x,y)\to(+\infty,+\infty)}\left(1+\frac{1}{xy}\right)^{x\sin y}.$\par
% p154-b26; source PDF page 154
\par\Needspace{5\baselineskip}
\blockanchor{p154-b26}
\majorquestion{二}计算题. 每题 10 分, 共 40 分.\par
% p155-b00; source PDF page 155
\blockanchor{p155-b00}
\question{1}计算定积分 $\int_0^{100\pi}\sqrt{1 - \cos 2x}\mathrm{d}x.$\par
% p155-b01; source PDF page 155
\blockanchor{p155-b01}
\question{2}计算 $\iint_{D} \sqrt{1 - \sin^2(x + y)} \, \mathrm{d}x \, \mathrm{d}y$, 其中 $D = \{(x, y) \mid 0 \leq y \leq x, 0 \leq x \leq \frac{\pi}{2}\}$.\par
% p155-b02; source PDF page 155
\blockanchor{p155-b02}
\question{3}已知  $\left\{\begin{aligned} x &= t^{2} + t, \\ y &= \ln(1 + t). \end{aligned}\right.$  求  $\frac{dy}{dx}$  与  $\frac{d^{2}y}{dx^{2}}$ .\par
% p155-b03; source PDF page 155
\blockanchor{p155-b03}
\question{4}已知 $u = f\left(xy, \frac{y}{x}\right)$, 其中 $f$ 二阶连续可微, 求 $\frac{\partial u}{\partial x}$ 与 $\frac{\partial^2 u}{\partial x \partial y}$.\par
% p155-b04; source PDF page 155
\par\Needspace{5\baselineskip}
\blockanchor{p155-b04}
\majorquestion{三}讨论题. 10 分.\par
% p155-b05; source PDF page 155
\par\Needspace{6\baselineskip}
\blockanchor{p155-b05}
设二元函数\par
\nopagebreak[4]
% p155-b06; source PDF page 155
\blockanchor{p155-b06}
\displayformula{f (x, y) = \left\{ \begin{array}{l l} (x ^ {2} + y ^ {2}) ^ {\alpha} \sin \frac {1}{x ^ {2} + y ^ {2}}, & x ^ {2} + y ^ {2} \neq 0; \\ 0, & x ^ {2} + y ^ {2} = 0. \end{array} \right.}
% p155-b07; source PDF page 155
\blockanchor{p155-b07}
分别讨论 $f(x,y)$ 在 $(0,0)$ 处连续，偏导数存在时，$\alpha$ 的取值范围.\par
% p155-b08; source PDF page 155
\par\Needspace{5\baselineskip}
\blockanchor{p155-b08}
\majorquestion{四}解答题. 每题 10 分, 共 40 分.\par
% p155-b09; source PDF page 155
\blockanchor{p155-b09}
\question{1}求 $\sum_{n=1}^{\infty} \frac{1}{n \cdot 2^n}$ 的和.\par
% p155-b10; source PDF page 155
\blockanchor{p155-b10}
\question{2}设 $f_{n}(x) = n^{\alpha}xe^{-nx}$，求 $\alpha$ 的取值范围，使得函数列 $\{f_n(x)\}$ 在[0,1]上一致收敛.\par
% p155-b11; source PDF page 155
\par\Needspace{6\baselineskip}
\blockanchor{p155-b11}
\question{3}判断级数或反常积分的敛散性.\par
\nopagebreak[4]
% p155-b12; source PDF page 155
\blockanchor{p155-b12}
\displayformula{\sum_ {n = 1} ^ {\infty} \frac {(- 1) ^ {n + 1}}{n} \frac {a}{1 + a ^ {n}} (a > 0). \tag {1}}
% p155-b13; source PDF page 155
\blockanchor{p155-b13}
\displayformula{(2) \int_ {0} ^ {\frac {1}{e}} \frac {1}{x ^ {p} \ln x} \mathrm{d} x (p > 0).}
% p155-b14; source PDF page 155
\par\Needspace{5\baselineskip}
\blockanchor{p155-b14}
\majorquestion{五}证明题. 10 分.\par
% p155-b15; source PDF page 155
\blockanchor{p155-b15}
已知函数 $f(x)$ 在 $(a,b)$ 内二阶可导，$f'(x)$ 在 $[a,b]$ 上连续，且满足 $f(a) = f(b) = 0$，$\int_{a}^{b} f(x) \, \mathrm{d}x = 0$，证明：存在 $\xi \in (a,b)$，使得 $f''(\xi) = f(\xi)$.\par
